\subsection{Indirect Support Generalizes Across Chromosomes and Recovers Disease Genes}

%1. Hold out cross validation to show indirect predicts direct
%2. Temporal to show combined predicts future
%3. Stratify across trait architectures and sample size, rare vs. common
%4. Bias towards well studied genes or traits with higher power
%5. Comparison to direct only baseline

% Figure 2 panels
%a. Show indirect predicts direct in hold one out chromosome (portal_model)
%b. Show mechanistic alignment in portal model results compared to null for combined
%c. Show relative success results compared to benchmarks for combined (portal model)
%d. Show it predicts future GWAS signal (portal) 

% Add how many more indirect genes vs only direct
% Plot - standard direct vs indirect plot and show orthongality.
% direct support is probablistic nearest gene
% Merge Cross chromosome lines for each trait
% Plot of sample size to demostrate calibration with direct support
% Another figure to show miami plot 

% What is a trait architecture?
% How do I show the bias towards well studied genes or traits with higher power?

% Think of threshold of cimbined as FDR (False Discovery Rate) based on Bayes Factors. TATTA paper has the FDR. 

% Relevant genes are if you perturb them they fuck up. However we can't perturn everything in humans to learn this because of ethics and cost.
% So how might relevant genes show up.
% - Successful drug targets would be a nice proxy, because they would target relevant genes no matter if they show up in GWAS.
% - Pertub seq data could also be leveraged maybe? That would support the causality of relevant gene definition. 


\begin{figure}[t!]
    \centering
    
\includegraphics[width=1\linewidth]{figures/placeholder-max-caption-2in.png}
    \caption{(a) A subset of pigean models run on a 13 trait common disease validation set. Displaying top performing models and models selected through assessment of top genesets. Marker size represents the complexity of the model and the multiple transparent grey boxes show overlayed null distrbutions in terms of both relative success and mechanistic alignment for each model (Methods). (b) Same models run over a 40 trait rare disease valiation set. (c) Over all model selection experiments the proportion of genesets that came from each geneset library that where significant for a given trait-model pair. }
    \label{fig:model-selection-main}
\end{figure}

% General overview paragraph of validation
We next sought to systematically validate PIGEAN along two goals: (i) predicting genes that are biologically and clinically meaningful and (ii) identifying gene sets that are mechanistically aligned with each phenotype. To do so, we evaluated a focused benchmark comprising $n=13$ common traits selected for representative clinical coverage and statistical power, and $n=40$ rare traits partitioned into $n=20$ disease-specific lists and $n=20$ disorder-category lists (see Methods). Because our full universe spans $>7,000$ traits, an exhaustive search over the model space was computationally infeasible; instead, we executed a targeted model-selection procedure that iteratively pruned candidates while answering specific research questions about PIGEAN’s behavior. For model selection, we defined a Minimum Description Length (MDL) criterion \cite{} that balances complexity and accuracy, where complexity captures both the number of gene sets made available to the model and the number of gene sets deemed significant for prediction, and accuracy integrates (a) potential clinical success of nominated targets and (b) mechanistic alignment between significant gene-set semantics and the phenotype under study. Operationally, this yielded three independent axes for comparison across models: (1) clinical success, quantified as the relative success \cite{minikel_refining_2024} of PIGEAN-prioritized genes that exceed a threshold of direct/indirect genetic support; (2) mechanistic alignment, quantified as the normalized discounted cumulative gain (nDCG) of semantic similarity ranks between significant gene-set names and their corresponding phenotype; and (3) parsimony, quantified as the balance between the number of predictive gene sets surpassing a $\beta$ threshold and the total size of the gene-set libraries included.

% Detailed description of top performing model with pointers to supplementary information. I think it helps ground the reader in the analysis when discussing just a single models performance
The top-performing model across the common and rare validation sets used two libraries: MSigDB \cite{} (with HPO-derived gene sets removed to avoid circularity) and 2024 Mouse Phenotypes \cite{}, for a total of $41,627$ gene sets. Under this configuration, the relative success for common disease was 8.623 (95\% CI: 5.499–13.534), representing a near 4-fold increase over prior reports \cite{minikel_refining_2024}, and for indirect-only genetic support it was 3.255 (95\% CI: 1.748–6.059), a regime that, to our knowledge, has not been previously characterized. Further, from rare traits using only a curated set of rare gene lists, we achieved a relative success of 2.167 (95\% CI: 1.028-4.57) which is only from indirect-support only as direct support in this context leverages the rare disease gene lists which would inherently be circular. Owing to PIGEAN’s regularization, these clinical-success gains were driven by an average of $124$ significant gene sets per common trait and $139$ per rare trait (Supplementary Information), amounting to approximately $0.3\%$ of the $41,627$ available gene sets. It should be noted that these relative success results are only amoung the validation traits in this analysis, and upon analyzing the full phenome of $>7000$ traits which is reported in the subsquent sections.

Mechanistic alignment, measured via nDCG (idealized ranking $=1$), was $0.79\pm 0.09$ , and the average phenotypic similarity of the top $k=10$ $\beta$-ranked gene sets was $0.202$ for both common and rare phenotypes. Finally, increase of phenotypic similarity among significant gene sets over a random null was statistically significant for common traits ($p=6.36\times 10^{-19}$) and for rare traits ($p=1.36\times 10^{-5}$) demostrating PIGEAN is capturing a genetic signal that is supported by literature. Figure~\ref{fig:model-selection-main} shows a subset of evaluated models that are of particular interest, and it is visually clear that the MsigDB + Mouse Phenotype model performs well overall balancing clinical relevance, mechanistic alignment, and model complexity. 

\begin{figure}[t!]
    \centering
    \includegraphics[width=1\linewidth]{figures/indirect-support/real/temp-figure-3.png}
    \caption{Relative success of drug targets launched post-2018 over baseline Magma supported genes using non-cicurlar models using gene sets curated from pre-2018 sources priorized (a) with direct genetic support and (b) indirect genetic supprt with secondary $y$-axis showing number of predicted targets for each genetic support threshold. Improvement in regression fit ($\Delta R^2$) of regressing Magma Z-scores from GWAS pre-2018 on GWAS post-2018 with (c) a direct genetic support filter and (d) an indirect genetic support filter (outlier traits removed).}
    \label{fig:validation-figure-2}
\end{figure}

% Back to the future validation
Our final validation step was designed to assess the potential impact of circularity between GWAS results and curated geneset libraries in predicting drug targets as well as in forecasting future GWAS significance. Because many published genesets are curated from GWAS signals, we sought to confirm that PIGEAN was not simply reinforcing known signals but was instead capable of highlighting novel associations. Similarly, we repeated our drug target relative success analysis using only targets that would not have been available prior to the publication of a given geneset. To this end, we created 10 additional models using gene set sources from MsigDB, Mouse Phenotypes, and MeSH Terms. For each library, we curated pre-2018 and post-2018 versions based on provenance information. We then applied our validation analyses across 67 traits with both pre- and post-2018 GWAS. For the relative success analyses, only drug targets launched after 2018 were included to avoid target circularity.  

As shown in Figure~\ref{fig:validation-figure-2}a-b, we observe consistent improvements in relative success as a function of both direct and indirect support, even after temporally excluding any form of GWAS or drug target information. Notably, the pre-2018 text-mining model, which explicitly removed all GWAS and drug MeSH terms (Methods), achieved the strongest performance. At a direct genetic support log-odds threshold of 4 ($P=98.2\%$ relevance), this model reached a relative success of 4.51 (95\% CI: 2.67–7.63) across 50 predicted targets, significantly above the Magma baseline ($p$-value $=2.5\times 10^{-6}$). Although somewhat weaker, indirect genetic support also improved relative success, reaching 3.04 (95\% CI: 1.93–4.80) at the same log-odds threshold with 108 predicted targets. This highlights an important distinction between Figure~\ref{fig:validation-figure-2}a and~\ref{fig:validation-figure-2}b: indirect support identifies additional targets not captured by direct genetic support alone, as reflected in the plotted trend lines showing the number of predicted targets as a function of the (in)direct support thresholds. We further analyze the complementarity of direct and indirect support in expanding drug target diversity in the following section.  

We next evaluated how well direct and indirect genetic support could predict future GWAS significance by regressing on post-2018 MAGMA Z-scores. When directly regressing Z-scores (Methods), we observed modest but statistically significant gains in model fit ($\Delta R^2$) when incorporating indirect support, with maximum improvements across all traits of rangeing from 0.0019 to 0.017 (Extended Data Figure XXXX). The proportion of traits with significant $\Delta R^2$ also varied by model (ranging from 62.9\% to 91.4\%). However, substantial noise remained in non-significant genes, where MAGMA Z-scores were nearly random, such that even baseline regressions of pre-2018 versus post-2018 Z-scores yielded a low average $R^2 = 0.14$ across 67 traits.  

To test whether genetic support could act as a filter to reduce noise and forcast genes that will be predictive of future GWAS, we regressed post-2018 MAGMA Z-scores on pre-2018 scores while progressively restricting to genes meeting increasing thresholds of direct or indirect genetic support (Figure~\ref{fig:validation-figure-2}c-d). Across all models, regression slopes remained positive for both direct ($\beta$ range: 0.014–0.038) and indirect ($\beta$ range: 0.050–0.058) support, with indirect support providing a greater increase in model fit. Importantly, filtering on stronger evidence not only increased average $\Delta R^2$ but also strengthened the concordance of pre- and post-2018 associations, consistent with the notion that PIGEAN prioritizes genes with reproducible relevance. Together, these temporal validations support that PIGEAN is not simply recapitulating latent GWAS or drug target information embedded in curated libraries, but instead enriches for signals that are predictive of future discoveries.  
